\section{Results}

\subsection{What rice-atlas is, on inspection}

\begin{table}[h]
\centering
\begin{tabular}{@{}lll@{}}
\toprule
& rice-atlas & arabidopsis-atlas (this work) \\
\midrule
Primary dataset cited & none found & \NumOrgansWithGeometryCitation{}/\NumOrgansModeled{} organs cite a real structural paper \\
Spaceflight data & none & OSD-120, OSD-314 (real OSDR counts) \\
License (GitHub-detected) & \texttt{NOASSERTION} & MIT (this repository's own) \\
FAIR metadata & none & \texttt{LICENSE}, \texttt{CITATION.cff}, \texttt{.zenodo.json} \\
GitHub Pages & not deployed & deployed \\
Independent review & none found & ABAI evidence-gate attestation on file \\
\bottomrule
\end{tabular}
\caption{Attributes compared directly by inspecting both repositories, not by
description.}
\end{table}

\subsection{Real spaceflight response, by organ}

Of the five organs modelled, one (root) has a directly tissue-matched real spaceflight
comparison: OSD-120 profiled root tissue specifically, at Day 13, under
\NumFlightSamplesOsdOneTwentyZero{} flight and \NumGroundSamplesOsdOneTwentyZero{} ground
control samples across \NumGenesOsdOneTwentyZero{} genes. The largest mean flight/ground
CPM ratio in that comparison is \TopFoldChangeOsdOneTwentyZero{} $\log_2$ units, for
\texttt{\TopGeneOsdOneTwentyZero{}}.

OSD-314 profiled whole seedlings, not a single organ, across the same
\NumGenesOsdThreeFourteen{} genes, comparing \NumMicrogSamplesOsdThreeFourteen{} real
microgravity samples against \NumGroundSamplesOsdThreeFourteen{} 1g ground samples; its
largest ratio is \TopFoldChangeOsdThreeFourteen{} $\log_2$ units, for
\texttt{\TopGeneOsdThreeFourteen{}}. Because it is not organ-specific, we report it in the
data provenance ledger rather than attach it to any single organ in the viewer, to avoid
implying a tissue-level claim the data does not support.

The remaining three organs (rosette leaf, inflorescence axis, flower/silique) have no
organ-matched spaceflight dataset wired into this version of the atlas; the viewer states
this plainly for each rather than substituting the whole-seedling numbers as if they were
organ-specific.

\subsection{Verification of the viewer itself}

The viewer was checked interactively, not only by static code review: in both a local
development server and the actual production build served at the live GitHub Pages URL,
clicking each organ was confirmed to open an information panel showing the correct TraVA
category, the correct geometry citation, and -- for root -- the real gene table described
above, with no console errors in either environment.
